\section{Task and Data}\label{sec:data}

\paragraph{Task.}
The input is the English description of a CVE.
The outputs are the eight CVSS v3.1 base metrics (attack vector, attack complexity, privileges required, user interaction, scope, and the confidentiality, integrity, and availability impacts), the base score, the severity tier (\textsc{Critical}, \textsc{High}, \textsc{Medium}, \textsc{Low}), and the CWE identifier.
Figure~\ref{fig:task} shows one example.

\paragraph{Source.}
NVD records come from the Fraunhofer FKIE mirror of the NVD JSON feeds~\citep{fkie2026nvdfeeds}, pinned to release \texttt{v2026.10.05-000016}.
NVD revises records over time, and a fixed release keeps every number in this paper reproducible.

\paragraph{Filtering.}
A CVE is kept if it has an English description of at least 20 characters, a CVSS v3.x base metric with a severity tier, and a specific CWE.
The placeholders \texttt{NVD-CWE-noinfo} and \texttt{NVD-CWE-Other} do not count.
When a CVE carries several CVSS entries, the primary score assigned by NIST is used.
This leaves \nAll{} CVEs from 2022 to 2024.
\nVThreeZero{} of them use CVSS v3.0 and the rest v3.1, and the base-score equations of the two versions agree on every one.

\paragraph{Temporal split.}
CVEs with 2022 and 2023 identifiers form a pool of \nTrainVal{} CVEs.
It is split 80/20 at random, stratified by severity, into \nTrain{} training and \nVal{} validation CVEs.
All \nTest{} usable CVEs with 2024 identifiers form the test set.
Validation data serve only to choose hyperparameters and checkpoints.
Table~\ref{tab:data} gives the label distribution of each split.

\begin{table}[t]
\centering
\caption{Data splits, with severity columns giving the share of CVEs (\%) in each tier.}
\label{tab:data}
\small
\begin{tabular}{lrrrrrr}
\toprule
Split & CVEs & Critical & High & Medium & Low & CWEs \\
\midrule
Train (2022--23) & 40{,}436 & 15.4 & 39.3 & 43.7 & 1.7 & 464 \\
Validation (2022--23) & 10{,}109 & 15.4 & 39.3 & 43.7 & 1.7 & 293 \\
Test (2024) & 36{,}325 & 13.2 & 36.2 & 48.6 & 1.9 & 512 \\
\bottomrule
\end{tabular}
\end{table}


\paragraph{Label distribution.}
\textsc{Low} is under 2\% of every split.
CWEs follow a long tail: the training set has \nTrainCwes{} distinct CWEs, the test set \nTestCwes{}, and the 100 most frequent training CWEs cover \nTopCov\% of test CVEs.

\paragraph{CVSS calculator.}
The CVSS v3.1 base-score equations~\citep{first2019cvss31} are implemented in full, including the floating-point-safe round-up that the specification requires.
Applied to the true vectors, the calculator reproduces the NVD base score and severity for all \nAll{} CVEs.
Any error in a formula-based score therefore comes from the predicted vector.

\paragraph{Excluded baseline.}
VLAI~\citep{bonhomme2025vlai} is a public severity classifier.
Its training data, the \texttt{CIRCL/vulnerability-scores} dataset, uses a random split over all years, and rows sampled from the training split include CVEs from 2024, 2025, and 2026.
Scoring it on the 2024 test set would measure memorization as much as generalization, so it is left out.
